\section{A source domain input and a sharp limitation}

Fix the primitive ratio, the large block constants and the variational
constants of \cite{B109}. The following result discharges the
common-radius version of the source-domain hypothesis used in the
homogeneous birth-radius argument.

\begin{proposition}[Uniform common-radius source inclusion]
There are fixed $c_*,a_*>0$ such that, with
$a=\min(\alpha_0,\alpha_1,a_*)$, source conditions (i)--(iii) of
\cite[pp.~262--263, (1.10)--(1.16)]{B109}, imposed with all three
parameters equal to $c_*a$, imply membership in the full space
$\mathcal U_j^c(X,\alpha_0,\alpha_1)$.
The constants are uniform in admissible $X$, refinement and finite
volume, and in the target radii below the fixed source ceilings.
\end{proposition}

\begin{proof}
The displayed estimate on p.~263 immediately after (1.16) applies the
source variational Proposition~9. Each reconstructed curvature and
current quantity in condition (iv), for both the complex field and
its real reference, is at most $2B_3\alpha'_0$ times its respective
mesh weight in (1.16). Here $B_3$ is fixed with the source geometry;
the estimate is used below its fixed smallness ceiling. Choose $a_*$
inside that ceiling and
$c_*\leq\min(1/4,(4B_3)^{-1})$. Conditions (i)--(iii) at the target
radii follow by monotonicity. All normalized reconstructed bounds
are at most $a/2\leq\alpha_0/2$, proving (iv). This choice is independent
of the shrinking radii. This argument uses the quantitative display,
not only the subsequent qualitative smaller-constant sentence.
\end{proof}

\begin{corollary}[An independent physical product ball]
Write $U=\exp(i\eta A)$ and use the real reference $I$, with $J$
independent of $A$. For a fixed $c>0$, the ball
\[
 \norm{A}_{\infty}+\norm{\nabla_\eta A}_{\infty}<ca,
 \qquad \norm{J}_{\infty}<ca
\]
lies in $\mathcal U_j^c(X,\alpha_0,\alpha_1)$. The norms, mesh factors
and current convention are those of the source; the fields are given
on the required collars. Straight segments in these independent
variables stay in the ball. In particular, the previously proved
homogeneous potential and nonlinear-current estimates imply that a
source function bounded by $A_X$ specializes with the same bound to
the full coarse-link product ball
$\norm{B}_\infty<c'\min(\alpha_0,\alpha_1,a_*)$.
\end{corollary}

\begin{proof}
The real reference has zero curvature. Plaquette multiplication gives
\[
 \norm{U(\partial p)-I}
 \leq C\eta^2
 \bigl(\norm{\nabla_\eta A}_\infty+\norm{A}_\infty^2\bigr).
\]
For $a_*\leq1$, a fixed small $c$ puts this and the independent current
strictly inside the smaller conditions (i)--(iii). Apply the proposition.
The product ball is convex in $(A,J)$. For the last assertion substitute
the existing uniform holomorphic bound
$\norm{A_{\rm ord}(B)}_{C^1,\eta}
+\norm{J_{\rm phys}(U_{\rm ord}(B))}_\infty\leq C_h\norm{B}_\infty$.
\end{proof}

This proves only the stated physical product-domain inclusion.
Locality of the specialized function in $B$, noisy-current bounds,
mixed-boundary geometries and all complex cluster-decoupling domains
are separate issues. Uniformity in a varying primitive ratio is not
asserted.

\begin{proposition}[Shrinking envelopes do not imply an $O(g)$ score]
The analytic bounds with
$\alpha_{a,j}=C_a g_j(\log g_j^{-2})^{q_a}$ of
\cite[(2.27)--(2.31)]{B119}, even with gauge invariance and vanishing
vacuum and Maxwell jets, do not alone imply an $O(g_j)$ fluctuation
score or an $O(g_j)$ interaction-jet norm through two background
derivatives.
\end{proposition}

\begin{proof}
It suffices to test the asserted implication at a fixed finite mesh.
For a plaquette holonomy $P$, set
\[
 Q(P)=\tau(2I-P-P^{-1}),\qquad
 E(P)=\frac{\lambda}{C_*\alpha_0^4\eta^8}Q(P)^2,
\]
where $\lambda\ne0$ is fixed and small. The identity
$P+P^{-1}-2I=(P-I)^2P^{-1}$ shows that $|E|\leq|\lambda|$ on the
source curvature domain after fixing $C_*$. This is holomorphic and
gauge invariant, is independent of $J$, and begins at fourth order
at the identity. A fixed localization domain supplies the prescribed
decay factor; identical copies over plaquettes supply Euclidean
covariance. Its Maxwell and vacuum extractions vanish.

On a commuting direction $P(t)=\exp(i\eta^2tH)$ with
$\tau(H^2)>0$, its leading term is
$c_H\lambda t^4/\alpha_0^4$, with $c_H>0$. For the physical
fluctuation line $t=B+gx$, let $V(B,x)=E(B+gx)-E(B)$. At $B=0$,
the leading terms of $V,\partial_xV,\partial_BV,\partial_B^2V$ are
respectively
\[
 \frac{c_H\lambda}{\alpha_0^4}
 \bigl(g^4x^4,\ 4g^4x^3,\ 4g^3x^3,\ 12g^2x^2\bigr).
\]
The higher terms have an extra $g^2$ on fixed $x$ intervals. With
$\alpha_0=C_0g(\log g^{-2})^{q_0}$, the first two are inverse powers
of $\log g^{-2}$, while the last two have additional factors $g^{-1}$
and $g^{-2}$. These are not $O(g)$, also in fixed local Gaussian
$L^p$ norms. The zero quadratic jet rules out a repair by Maxwell
subtraction alone.
\end{proof}

The example disproves an inference from the stated function class;
it is not claimed to be generated by the Wilson initial state.
For an actual chart and a bounded analytic insertion, elementary
Cauchy estimates yield costs $g/r$ and $g^2/r^2+g^2/r$ for the first
two fluctuation derivatives on a fractional interior. At
$r\asymp g(\log g^{-2})^q$ these are inverse logarithms. Weighted
local Hessian bounds additionally require support summation. The
fixed-radius reasoning in \cite[(1.19)--(1.22), Lemma~1]{B116}
cannot silently make these shrinking-domain constants uniform.

\begin{proposition}[An exact normalized one-mode diagnostic]
Convexity and locality, together with a bounded shrinking analytic
interaction of zero quadratic jet, do not imply an absolute $O(g^2)$
static-subtracted response Hessian for a mark of fixed analytic norm.
\end{proposition}

\begin{proof}
This is a one-mode test, not a Wilson-generated state. Fix $\lambda>0$,
set $W(u)=u^2$, and use the normalized measure
\[
 d\mu_B(x)=Z_B^{-1}
   e^{-x^2/2-\lambda(B+gx)^4/r^4}\,dx,
 \qquad R(B)=\mathbb E_{\mu_B}\bigl[(B+gx)^2-B^2\bigr].
\]
Its real fluctuation Hessian is
$1+12\lambda g^2(B+gx)^2/r^4\geq1$. The interaction is bounded
by $\lambda$ on the physical complex ball $|u|<r$ and has zero
quadratic jet; $W$ has a fixed norm on a fixed ball. Put
$\epsilon=\lambda g^4/r^4$ and $m_n=\mathbb E_{\mu_0}x^n$.
Normalized differentiation gives
\[
 \left.\partial_B\mathbb E_{\mu_B}x\right|_0
       =-\frac{4\epsilon}{g}m_4,
 \qquad
 \left.\partial_B^2\mathbb E_{\mu_B}x^2\right|_0
       =-\frac{12\epsilon}{g^2}(m_4-m_2^2)
        +\frac{16\epsilon^2}{g^2}(m_8-m_2m_6).
\]
Using $R(B)=2Bg\mathbb E_{\mu_B}x+g^2\mathbb E_{\mu_B}x^2$,
\[
 R''(0)=-\epsilon(28m_4-12m_2^2)
          +16\epsilon^2(m_8-m_2m_6)
        =-72\epsilon+O(\epsilon^2).
\]
The final equality follows by dominated Gaussian integration as
$\epsilon\downarrow0$. For $r=g(\log g^{-2})^q$, this is an inverse
logarithm and is not $O(g^2)$.
\end{proof}

This diagnoses the missing interaction-jet input; it does not
contradict the response theorem, whose $O(g)$ jet hypothesis fails
in this example. Its displayed sign is not an amplification claim.

There is a quantitative reason to retain this distinction. Along
$\gamma_{k+1}=\gamma_k-1$ down to a fixed weak endpoint, a product of
allowed bounds $1+C/(\log\gamma_k)^a$ grows at least as
$\exp[c\gamma_i/(\log\gamma_i)^a]$. No birth suppression
$\exp[-c'(\log\gamma_i)^\nu]$ with fixed finite $\nu$ pays that
majorant. A verified stronger norm after a finite startup, or a
$g^2$ times polylogarithm propagation bound, is still required.

